% Transcription — fonds Alexandre Grothendieck, shelfmark 29, batch 3 (pages 41-60).
% Established from the facsimile archives/batches/29/batch-03.pdf (a mirror of
% https://grothendieck.umontpellier.fr/29.pdf), per the skill
% transcribe-grothendieck.
%
% Fourteen of the twenty pages carry mathematics. Pages 48, 50, 52, 54, 56 and
% 58 are six consecutive sheets of an unrelated violet typescript — "n° 385",
% pages 10, 11, 12, 13, 22 and 23, on distributions on manifolds of class C^r —
% and Grothendieck used their versos for pages 47, 49, 51, 53, 55 and 57. They
% carry nothing in his hand, so they are not transcribed and their numbers
% simply do not appear.
%
% The batch falls into three runs.
%
% Pages 41 to 45 continue the sheaf "Finitude (Pbs et conjectures)" begun at
% page 39 of batch 2. Page 42 finishes the list of conjectures started on page
% 40 — its item c) is the birational-invariance conjecture for tame
% ramification — and opens "2) Groupe fondamental local"; pages 43 and 45,
% each struck through with a single long pencil stroke, construct the two
% quotient categories Coh(X/X_0) and R(X/X_0) of a formal scheme modulo a
% scheme of definition, and compute both in the proper case by the
% comparison theorem; page 44 states "3) Théorèmes de finitude et
% simultanéité", the constructibility of pi_1 of the fibres. Page 41 is a
% scratch sheet on the completions of a ring in two variables, later crossed
% with heavy pencil doodling.
%
% Pages 46 to 57 are the run Grothendieck headed "Th. de finitude relatif pour
% pi_1" on page 46 and paginated himself I to VI: Lemme 1 (descent of an
% isomorphism of the geometric generic fibres to a finite étale cover of the
% base, with SGA IX 6.1-6.2 invoked and SGA X expressly set apart), Lemme 2
% (constructibility of f_*), Lemme 3 (H^1 along a finite partition into locally
% closed parts), the four conditions (i) to (iv) of a "Main Proposition", and
% its two corollaries. The hand is fast and heavily struck, and the scan is
% 228 ppi; a good deal of it is left as \ill{}, which is the honest reading and
% not a shortcut.
%
% Pages 59 and 60 open a new sheaf, "pi_1 et Pic^tau birationnels", of which
% only Grothendieck's own page 1 is in this batch: pi_1' as a projective limit
% of the discrete quotients of order prime to p, reduced by van Kampen to the
% affine normal case and then by a generic linear section to a curve.
%
% Two English words stand in the French — "wild" for wild inertia on page 60,
% and "Main Proposition" on page 51 — and are left as written, as is the
% "kif kif" of Lemme 3.
%
% The dating is the inventory's own, for the whole shelfmark.
%
% Pass: Opus 5 (claude-opus-5), 2026-08-22 — first pass, unchecked against the pages by a human.
\documentclass[11pt,a4paper]{article}
% Le préambule commun à toutes les transcriptions du fonds Grothendieck.
%
% Il tient en une idée : **l'incertitude se compose, elle ne se résout pas.**
% Une transcription qui lisse un mot douteux a détruit la seule chose pour
% laquelle on la faisait — un lecteur doit pouvoir savoir, à chaque endroit,
% ce qui était sur le feuillet et ce qui a été ajouté. Les six macros qui
% suivent sont donc l'appareil critique, et non de la décoration.
%
% Les mêmes commandes sont reconnues par `scripts/render.mjs`, qui produit la
% vue de lecture du site. Une macro ajoutée ici doit l'être aussi là-bas,
% faute de quoi le rendu échouera bruyamment — ce qui est voulu : un rendu
% silencieusement faux est pire qu'un rendu absent.

\ProvidesPackage{grothendieck}[2026/08/08 Transcriptions du fonds Grothendieck]

\usepackage{amsmath,amssymb,amsthm}
\usepackage{mathrsfs}
\usepackage{stmaryrd}
% Les diagrammes commutatifs sont partout dans ce fonds : c'est la langue
% même de la géométrie algébrique de Grothendieck, et une transcription qui
% les rendrait en prose ne transcrirait rien.
\usepackage{tikz-cd}
% Français par défaut — c'est la langue des feuillets — mais l'anglais est
% chargé aussi : la traduction et le résumé sont des documents anglais, et la
% typographie française (espaces avant les deux-points) y serait une faute.
% Ces documents basculent par \selectlanguage{english} en tête de corps.
\usepackage[english,french]{babel}
\usepackage{fontspec}
\usepackage{geometry}
\geometry{a4paper,margin=25mm,marginparwidth=28mm}
\usepackage{marginnote}
\usepackage{xcolor}
% ulem, not soul. `soul`'s \st and \ul are built on a character-by-character
% scan that XeLaTeX's font machinery breaks: the document fails with an
% undefined control sequence rather than degrading. ulem does the same two jobs
% and is Unicode-engine safe. `normalem` stops it hijacking \emph, which the
% transcriptions use for Grothendieck's own emphasis.
\usepackage[normalem]{ulem}
\usepackage{hyperref}
\hypersetup{colorlinks=true,linkcolor=black,urlcolor=black,pdfborder={0 0 0}}

% --- Métadonnées du lot -----------------------------------------------------
% Reprises telles quelles par le script de rendu, qui les affiche en tête de
% la vue de lecture. Elles fixent ce que le fichier prétend couvrir : sans
% elles, une transcription flotte sans point d'attache dans un fonds de seize
% mille feuillets.

\newcommand{\folder}[1]{\gdef\grothFolder{#1}}
\newcommand{\batch}[1]{\gdef\grothBatch{#1}}
\newcommand{\leaves}[2]{\gdef\grothFirst{#1}\gdef\grothLast{#2}}
\newcommand{\foldertitle}[1]{\gdef\grothFoldertitle{#1}}
\gdef\grothFolder{?}\gdef\grothBatch{?}\gdef\grothFirst{?}\gdef\grothLast{?}\gdef\grothFoldertitle{}

% --- Le résumé --------------------------------------------------------------
%
% Il ouvre la lecture modernisée, sous le titre « Résumé », et il est composé
% en petit. Ce n'est pas un ornement : c'est le seul passage du document qui
% ne soit pas des mathématiques, et un lecteur doit pouvoir le reconnaître
% avant de l'avoir lu — puis le sauter s'il connaît déjà le sujet, ou s'y
% arrêter s'il ne le connaît pas. Le filet qui le suit marque la reprise du
% corps.

\newenvironment{resume}
  {\small\setlength{\parskip}{0.45em}}
  {\par\medskip\hrule\medskip}

% --- Mots-clés ---------------------------------------------------------------
%
% En anglais, en clôture du résumé de la lecture modernisée : le vocabulaire
% moderne sous lequel chercher ce que les feuillets construisent. Le manifeste
% du site (`npm run manifest`) les extrait pour étiqueter la cote entière —
% cette ligne est la source unique des tags, il n'y a pas de fichier à côté.
\newcommand{\keywords}[1]{\par\smallskip\noindent
  {\footnotesize\textsc{Keywords} --- \textit{#1}}\par}

% --- L'appareil critique ----------------------------------------------------

% Le numéro de feuillet, en marge. C'est l'ancre qui permet de régler un
% désaccord sur une formule en regardant *un* feuillet plutôt qu'une chemise
% de six cents. Le site s'en sert aussi pour tourner les pages du fac-similé
% quand on fait défiler la transcription.
\newcommand{\leaf}[1]{%
  \marginnote{\footnotesize\color{gray!70}\textsf{#1}}%
  \label{leaf:#1}\ignorespaces}

% Illisible : on ne devine pas. Un mot inventé qui se lit comme les autres est
% le pire résultat possible.
\newcommand{\ill}{\textcolor{red!60!black}{[\,\ldots\,]}}

% Lecture douteuse : le mot est donné, et signalé comme incertain.
\newcommand{\uncertain}[1]{\uline{#1}}

% Ajout de l'éditeur — une lettre manquante, un mot sous-entendu.
\newcommand{\add}[1]{\textcolor{blue!50!black}{[#1]}}

% Biffé par Grothendieck lui-même. Ce qu'il a rayé fait partie du document :
% une rature dit souvent où la pensée a changé de direction.
\newcommand{\struck}[1]{\sout{#1}}

% Note du transcripteur, sur l'état matériel du feuillet.
\newcommand{\note}[1]{\par\smallskip{\small\color{gray!60!black}\textit{[#1]}}\par\smallskip}

% Note marginale de Grothendieck, distincte du corps du texte.
\newcommand{\marginal}[1]{\par\smallskip\noindent\hspace{1em}%
  {\small\color{gray!60!black}#1}\par\smallskip}

% --- Titre ------------------------------------------------------------------

\AtBeginDocument{%
  \begin{center}
    {\large\bfseries Fonds Alexandre Grothendieck --- cote n\textsuperscript{o}~\grothFolder}\\[2pt]
    {\itshape\grothFoldertitle}\\[4pt]
    {\small feuillets \grothFirst--\grothLast\ (lot \grothBatch)}\\[2pt]
    {\footnotesize\color{gray!60!black}Transcription établie d'après le fac-similé numérisé
     par l'Université de Montpellier.\\
     Les lectures incertaines sont soulignées, les illisibles notées [\,\ldots\,],
     les ajouts entre crochets.}
  \end{center}
  \vspace{1em}}

\endinput


\folder{29}
\batch{3}
\pages{41}{60}
\dating{1959-1969}
\watermark{Édition de démonstration}
\foldertitle{Groupe fondamental [Autour de SGA 4 et SGA 7] : lettres (1959, 1967, 1969), tapuscrits et copies de tapuscrit annotés (s.d.), notes manuscrites (s.d.)}

\begin{document}

\section*{« Finitude (Pbs et conjectures) », suite (pages 41 à 45)}

\page{41}
\note{la page est traversée de traits de crayon gras et d'un griffonnage qui
recouvrent une partie de l'écriture ; ce qui suit est ce qui en reste lisible}

$A$ anneau noethérien, $x, y \in A$

a) $A/x$, $A/y$ réguliers intègres de dim 1

b) $A$ \struck{complet} \ill{} pour la topologie définie par $xy$
\note{un mot est porté en interligne au-dessus du mot rayé, et n'est pas
lisible : c'est l'hypothèse elle-même qui reste indéterminée}

\begin{tikzcd}
\hat{A}^{(x)} & & \hat{A}^{(y)} \\
& A \arrow[ul] \arrow[ur] &
\end{tikzcd}

\[ \hat{A}^{(x,y)} = k[[x,y]] \]

$\hat{A}^{x} = \left\{ \sum a_{\alpha\beta}\, x^{\alpha} y^{\beta} \right\} = k[y][[x]]$
(seul un nb fini de termes, pour $\alpha \leq N$ donné, sont non nuls)

$A = \left\{ \sum a_{\alpha\beta}\, x^{\alpha} y^{\beta} \right\}$
(seul un nb fini des $a_{\alpha\beta}$, pour $\alpha + \beta \leq N$ donné,
sont non nuls)

$\hat{A}^{y} = \left\{ \sum a_{\alpha\beta}\, x^{\alpha} y^{\beta} \right\} = k[x][[y]]$
(seul un nb fini des termes, pour $\beta \leq N$ donné, sont non nuls)

\note{à droite de la page, un croquis à l'encre : deux axes portés $\beta$ et
$\alpha$, et une courbe qui découpe la région des indices}

\page{42}
\marginal{Prouvé}
c) \uncertain{Modérato} \struck{conjecture} d'invariance birationnelle de la
ramification modérée [si $X$ non \struck{\ill{}} projectif] (qui donne un sens
au pb.) et procéder comme b), a) [\struck{Lefschetz, résolution}
normalisation, Bertini\ldots] on est ramené \ill{} à une \struck{surface}
\emph{courbe} \struck{\ill{}}, \struck{diviseur à croisements normaux}, ok.

2) \emph{Groupe fondamental local.}

$X$ \struck{spectre} schéma noeth., \{ complet ; str.\ local et \emph{bon} \}.
$U$ un ouvert. \ill{}

\begin{itemize}
\item[a)] $\pi^t_1(U)$ de type fini
\item[b)] $\pi'_1(U)$ de présentation finie
\end{itemize}

Via résolution, et conjecture de \uncertain{fameuse} (invariance
birationnelle) dans le cas a), on doit s'en tirer, en se ramenant aux énoncés
analogues pour le 1)\ldots

\marginal{vérifier si c'est \emph{nécessaire}}

a) Par descente, on est ramené au cas $X$ normal, et par
\uncertain{hyperplan} [du moins si $U = X - \text{origine}$] on \ill{} où
$\dim X = 2$. On dispose alors de \uncertain{Abhyankar} et cela doit
marcher \ldots

\page{43}
\note{la page entière est barrée d'un long trait de crayon}

$\mathfrak{X}$ \add{pré}schéma formel loc.\ noeth., $\mathfrak{X}_0$ un
sous-schéma de définition, $\mathfrak{X}_0 = V(\mathcal{J})$.

On va définir la catégorie $\mathrm{Coh}(\mathfrak{X}/\mathfrak{X}_0)$ des
« Modules cohérents finis de $\mathfrak{X}$ mod $\mathfrak{X}_0$ ». Soit
$\mathrm{Coh}(\mathfrak{X})$ la catégorie des \struck{schémas formels finis}
Modules cohérents sur $\mathfrak{X}$. \struck{Soit} Soit $\mathfrak{S}$
l'ensemble des flèches $u$ de $\mathrm{Coh}(\mathfrak{X})$ telles que
$\operatorname{Ker} u$ et $\operatorname{Coker} u$ soient \struck{\ill{}}
\ill{} tués par une puissance de $\mathcal{J}$. Alors l'\uncertain{ens.}
$\mathfrak{S}$ satisfait aux axiomes qu'il faut pour former une catégorie de
fractions [N.B.\ on \ill{} ; en pratique on prend la sous-catégorie épaisse
des Modules cohérents tués par une puiss.\ de $\mathcal{J}$].
$\mathrm{Coh}(\mathfrak{X}/\mathfrak{X}_0)$ est la catégorie des fractions.

\emph{Exemple.} Soit $X$ propre sur $Y = \operatorname{Spec} A$, $A$ noeth.\
séparé et complet pour l'idéal $I$, $Y_0 = V(I)$, $X_0 = X \times_Y Y_0$,
$\mathfrak{X}$ le complété formel de $X$. Alors

\begin{tikzcd}
\mathrm{Coh}(\mathfrak{X}) & \mathrm{Coh}(X) \arrow[l, "\approx"'] \\
\mathrm{Coh}(\mathfrak{X})_{\mathrm{nilp}} \arrow[u] & \mathrm{Coh}(X)_{\mathrm{nilp}} \arrow[u]
\end{tikzcd}

par le th.\ d'existence, $\mathrm{Coh}(X)_{\mathrm{nilp}}$ étant les Modules à
support $\subset X_0$. On trouve
$\mathrm{Coh}(\mathfrak{X}/\mathfrak{X}_0) \approx \mathrm{Coh}(X - X_0)$ !

\page{44}
3) \emph{Théorèmes de finitude et \uncertain{simultanéité}}

Comment les groupes fondamentaux des fibres (d'un morphisme de type fini :
on désire des résultats genre constructibilité), qu'on \ill{} a), b), c) :

\marginal{c'est fait}

\begin{itemize}
\item[a)] Majoration \struck{« group-theoretical »} \emph{uniforme} dans les
cas de finitude \struck{dans} $\pi_1(X_{\bar s})$, si $X/S$ propre O.K.

\item[b)] Majoration \emph{uniforme} \struck{dans} des cas de finitude et de
\ill{} utiles pour les $\pi'_1(X_{\bar s})$, ou \emph{mieux} : la structure de
$\pi'_1(X_{\bar s})$ est fonction constructible de $s$

\item[c)] Majoration uniforme dans les cas de finitude de
$\pi^t_1(X_{\bar s})$.
\end{itemize}

Il semble que les méthodes qui résoudront 1) résoudront également 3).

\page{45}
\note{la page entière est barrée d'un long trait de crayon}

On définit une catégorie $R(\mathfrak{X}/\mathfrak{X}_0)$ ainsi.
$R(\mathfrak{X}) = $ catégorie des $\mathfrak{X}$-préschémas formels finis sur
$\mathfrak{X}$. On définit la catégorie des flèches $\mathfrak{S}$ [isom.\ mod
$\mathfrak{X}_0$] par : $u : \mathfrak{X}' \to \mathfrak{X}''$ est un isom.\
mod $\mathfrak{X}_0$ sss l'homom.\ d'Algèbres
$\mathcal{A}'' \to \mathcal{A}'$ correspondant est un isom.\ de Modules mod
$\mathfrak{X}_0$. On vérifie qu'on obtient un calcul de fractions à droite
[\uncertain{détail} bien compliqué, \uncertain{détail} bien \ill{} de bien,
explication]. La catégorie quotient est notée
$R(\mathfrak{X}/\mathfrak{X}_0)$.

\emph{Exemple.} Reprenons l'exemple précédent. Alors
$R(\mathfrak{X}) \approx R(X)$ par le Théorème de comparaison, donc
$R(\mathfrak{X}/\mathfrak{X}_0)$ est équivalente à $R(X/X_0)$, qui à son tour
est équivalent à $R(X - X_0)$.

\marginal{à vérifier}

[$R(X/X_0) \to R(X - X_0)$, c'est \ill{} \emph{facile} comme on vérifie
directement, c'est \ill{} comme on \uncertain{voit} \ill{} : le \ill{} global
du Main Theorem par ex.\ \ldots]

\section*{« Th. de finitude relatif pour $\pi_1$ » (pages 46 à 57)}

\page{46}
Th.\ de finitude \emph{relatif} pour $\pi_1$

\page{47}
\note{Grothendieck numérote cette page I, et les suivantes II à VI}

\struck{\ill{}}

\marginal{fid.\ plat}

\emph{Lemme 1.} Soit $f : X \to Y$ un morphisme \struck{\ill{}} de
préschémas loc.\ noethériens, \struck{de type f.},
\struck{à fibres géom.\ connexes}, \ill{} noeth.\ avec $Y$ intègre de point
gén.\ $y$, et \struck{soient} soient $X'$, $X''$ deux \uncertain{sous-schémas}
de $X$, tels que $X'_{\bar y}$ et $X''_{\bar y}$ soient
$K_{\bar y}$-isomorphes \ill{}. \struck{Alors $\exists$ revêtement fini étale
\ill{}} $Y_1$ de $Y$ tel que $X' \times_Y Y_1$ et $X'' \times_Y Y_1$ soient
$X_{Y_1}$-isomorphes.

Plus précisément, \ill{} tel homom.\ donné
$\bar u : X'_{\bar y} \to X''_{\bar y}$, si $K_1$ est la plus petite
\struck{sous-extension} de $k(\bar y)$ sur laquelle il est
\uncertain{défini}, \ill{} alors \struck{si} $K_1$ est \emph{étale} sur $Y$,
\ill{} et $Y_1$ \ill{} tel que $\bar u$ provient \ill{}
$X' \times_Y Y_1 \to X'' \times_Y Y_1$.

\ill{} d'un « propre » \ill{}

\emph{P.} --- \ill{} le cas, en introduisant $\underline{\mathrm{Hom}}_X(X', X'')$
\struck{$\underline{\mathrm{Isom}}_X(X', X'')$}, on peut supposer que
$X' = X$, d'autre part, on peut supposer inversement $Z = X''$ comme \ill{}
(quitte à le vérifier pour \ill{} comme \ill{} de l'image de
$X_{\bar y} \xrightarrow{\ \bar u\ } Z_{\bar y} \to Z$). L'hyp.\ indique que
(SGA IX 6.2.) $Z_y$ \ill{} est la fibre \ill{}. \struck{$Z$}
$X \times_Y Y_1$ \ill{} peut supposer contenue dans $k(\bar y)$, \ill{} que
\ill{} soit engendré par \ill{} $k(y) \to k$ \ill{}

\page{49}
Soit $Y_1$ le revêtement de $Y$ en $y_1$, \ill{}

\begin{tikzcd}
X \arrow[d] & X_1 \arrow[l] \arrow[d] \\
Y & Y_1 \arrow[l]
\end{tikzcd}

$Z \approx X \times_Y Y_1$, $Y_1$ \ill{}. \ill{} $X_1$ est étale sur $X$
propre, et les \struck{\ill{}} \ill{} de profondeur \ill{} donnent un
$X$-isom.\ $X_1 \approx Z$, ce qui achève la démonstration.

\marginal{\struck{\ill{}}}

NB.\ \ill{} on \ill{} \ill{} :

\[ \pi_1(X_{\bar y}, \xi) \to \pi_1(X, \xi) \to \pi_1(Y, \zeta) \to 0 \]

\struck{Cela} généralise \ill{} SGA IX 6.1., 6.2.\ \ldots Ne pas confondre
avec les résultats de SGA X, où on prend une fibre \struck{générique}
\uncertain{vraie} (en \ill{} il y faut supposer que $f$ \ill{} propre
\ldots).

\page{51}
\emph{Lemme 2.} Soit $X \to Y$ un morphisme de type fini, avec $Y$ loc.\
noeth.\ et univ.\ japonais. Alors pour tout faisceau d'ensembles constructible
$F$ sur $X$, $f_{*}(F)$ est constructible.

\marginal{Standard}

\marginal{univ.\ japonais}

\emph{Lemme 3.} $X$ schéma connexe \struck{\ill{}} noeth., $(X_i)$ une
partition finie en parties loc.\ fermées, $G$ un gpe fini. \struck{Supposons que} \struck{Considérons
l'ensemble des} \struck{$H^1(X_i, G)$ sont finis.} \struck{Pour tout $i$,}
Alors \struck{\ill{}} l'application

\[ H^1(X, G) \longrightarrow \prod_i H^1(X_i, G) \]

est à \ill{} : fibres \emph{finies}. En particulier, si pour tout $i$,
$H^1(X_i, G)$ est fini, kif kif pour $H^1(X, G)$.
\note{« kif kif » est bien ce que porte la page}

\ill{} Théorie de recollement des faisceaux étales, en utilisant lemme 2. Et
le fait suivant : \struck{Lemme 4.} Si $F$, $G$ sont des faisceaux étales
constr.\ sur $X$ noethérien, alors $\underline{\mathrm{Hom}}(F, G)$ est un
\ill{} \uncertain{ens.} fini.

\struck{Proposition} \struck{V\textsuperscript{o} Main Proposition} $S$ schéma
\ill{} noeth.\ univ.\ japonais, \ill{} $G$ un gpe fini. Considérons \ill{}

\page{53}
\begin{itemize}
\item[(i)] Pour tout $X$ de t.f.\ sur $S$, $H^1(X, G)$ est fini.

\item[(ii)] Pour tout sous-schéma \struck{\ill{}} \uncertain{intègre affine}
$S_0$ de $S$, $\ne \emptyset$ \ill{} tel \ill{} $S'$ de $S_0$, tout schéma
\add{affine} lisse de t.\ fini $X$ sur $S'$, à fibres géom.\ connexes non
vides, $H^1(X, G)$ est fini.

\item[(iii)] Pour tout $S'$, $X$ comme dessus, $S'$ \ill{} et tout
\uncertain{ouvert} non vide $U$ de $X$, il existe un \ill{} $U'$ de $U$ tel
que $H^1(X|U', G)$ soit fini.
\end{itemize}

\uncertain{Je} \ill{} bien ceci qui \ill{}

\begin{itemize}
\item[(iv) a)] Pour tout $s \in S$, et tout \uncertain{schéma} $\bar X$ de
\ill{} type fini non \uncertain{affine} connexe sur \ill{},
$H^1(\bar X, G)$ est fini

\item[b)] Pour tout \struck{\ill{}} schéma affine intègre $S'$
\struck{\ill{}} \add{comme dans (ii)}, et tout \ill{} $H$ \ill{},
$H^1(S', H)$ est fini.
\end{itemize}

\uncertain{Lorsque} $H$ est commutatif, on \ill{} affaiblir b) en y remplaçant
$H$ par $G$.

\emph{Dém.} \ill{} (ii) et (iii) résultent \uncertain{aussitôt} \ill{} lemme
3. Pour prouver que \ill{} au contraire pour (iv), on doit utiliser \ill{} et
plus lemme 1.

\page{55}
\emph{Cor\textsuperscript{1}} Supposons que $G$ est premier aux car.\
résiduelles. \struck{Alors} \struck{\ill{}} Supposons que pour tout
sous-schéma affine intègre non \ill{} $S_0$ de $S$, et tout \emph{ouvert}
\uncertain{dense} $S'$ de $S_0$, et pour tout \struck{\ill{}} sous-groupe $H$
de $G \times G$, on ait $H^1(S', H)$ fini. Alors pour tout $X$ de t.f.\ sur
$S$, $H^1(X, G)$ \ill{}. Or comme \ill{} ci-dessus, il suffit de faire \ill{}
la condition (iv) \ill{} $H$ \ill{}

\emph{Corollaire 2} Conditions équivalentes :

\begin{itemize}
\item[(i)] Pour tout $X$ de t.f.\ sur $Y$, et tout \uncertain{gpe} $G$ premier
aux car.\ résiduelles \ill{}, on a $H^1(X, G)$ fini

\item[(ii)] \struck{\ill{}} \uncertain{Idem}, avec $X$ \ill{} \ill{} affine
intègre $S_0$ et $S'_0$ \ill{}
\end{itemize}

NB.\ Dans le cor.\ 1, \ill{} on utilise le fait \ill{} pour \ill{} $h$ \ill{}
avec rg.\ de \ill{}, \ill{} pour tout $X$ de t.f.\ sur \ill{} et $G$ premier :
le cas de \ill{}, $H^1(X, G)$ est \ill{}

Voyons comment ceci peut être prouvé par les méthodes envisagées ici.

\page{57}
\marginal{\uncertain{finir} si le \ill{}}

\struck{\ill{}} Récurrence sur $\dim X$. \struck{\ill{}} \ill{} lemme 3, on
est \uncertain{ramené} au cas où \ill{} connexe. Si $\dim X = 1$, $X$ est
\uncertain{lisse} et \struck{\ill{}} \ill{} \struck{projectif} c'est facile
(réduction au cas de Abhyankar) \ill{} ayant un lemme de Serre--Lang.
\struck{\ill{}} c'est dû à \ill{} Serre--Lang.

Si $\dim X \geq 2$, on \ill{} une \uncertain{fibration} \emph{lisse}
$X \to Y$, $\dim Y < \dim X$, \ill{} si on \ill{} \ill{}, l'idée, de \ill{}
les \ill{} par \uncertain{hyp.\ de récurrence} \ill{}

\uncertain{Rmq.} \ill{} $X$ \ill{} sur $k$, \ill{} pour tout \struck{\ill{}}
\ill{} $X$ de \ill{} le cas de \ill{}. Alors la \emph{même} chose \ill{} fini.
\ill{} \ill{} tout $X$ de t.f.\ sur $k$ \ill{} est vrai pour \ill{}

\section*{« $\pi_1$ et $\mathrm{Pic}^\tau$ birationnels » (pages 59 et 60)}

\page{59}
$\pi_1$ et \emph{Pic}$^\tau$ birationnels

\page{60}
\note{Grothendieck numérote cette page 1}

\emph{Th.} Soit $X$ préschéma alg.\ sur $k$ alg.\ clos. Alors $\pi'_1(X) = $
limite projective des groupes \uncertain{quotients} discrets de $\pi_1(X)$
d'ordre premier à $p$, \ill{} est \ill{} de type fini.

Grâce à « van Kampen », on peut supposer $X$ affine, puis $X$ affine et
\emph{normal} [sinon, \ill{}
$X' \times_X X' \rightrightarrows X' \to X$, et $\pi_1(X)$ est engendré par
$\pi'_1(X')$ et un \ill{} fini de générateurs correspondant aux composantes
irréd.\ de $X' \times_X X'$, donc il en est de \ill{} \ill{} de type profini,
\ill{} $(p)$].

Si $\dim X = 0$, trivial.

Si $\dim X = 1$, théorie de la spécialisation du groupe fondamental nous
ramène au cas de la caractéristique 0, où on \struck{\ill{}} \ill{} les
résultats par voie transcendante.

Si $\dim X \geq 2$, on considère un revêtement \ill{}
$\bar X^n \longrightarrow \mathbf{P}^r_k$ et un \ill{} \struck{\ill{}}
\add{sous-espace lin.\ $L^{r-n+1}$} $H$, défini sur une ext.\ \emph{alg.\
close} $K$ de $k$, on peut \ill{}

\[ C = \bar X^n \times_{\mathbf{P}^r_k} L^{r-n+1} \]

qui \ill{} une courbe normale de \ill{} sur $K$, \ill{} « générique », et
alors on \ill{} que

\[ \pi_1(C) \longrightarrow \pi_1(X) \]

est \emph{surjectif} [\ill{}]. Donc $\pi'_1(C) \to \pi'_1(X)$ l'\ill{} aussi,
O.K.

\emph{Question} $X$ normale complète, $Y \subset X$ \struck{\ill{}} : alors
$\pi_1(X - Y)$ est-il le \ill{} groupe \ill{} \ill{} engendré \ill{} par les
groupes d'inertie \emph{wild} \ill{} et les composantes de \ill{} codim.\ 1 de
$Y$, est-il \ill{} top.\ de type fini ??

\end{document}
